\documentclass[10pt,a4paper]{article}

\usepackage[a4paper,margin=0.55in]{geometry}
\usepackage[T1]{fontenc}
\usepackage[utf8]{inputenc}
\usepackage{lmodern}
\usepackage[hidelinks]{hyperref}
\usepackage{enumitem}
\usepackage{titlesec}
\usepackage{parskip}
\usepackage{xcolor}

\setlength{\parindent}{0pt}
\setlength{\parskip}{0.25em}
\setlist[itemize]{leftmargin=1.2em, itemsep=0.12em, topsep=0.12em, parsep=0pt, partopsep=0pt}
\titleformat{\section}{\large\bfseries}{}{0pt}{}[\titlerule]
\titlespacing*{\section}{0pt}{0.6em}{0.35em}
\linespread{0.95}
\hypersetup{colorlinks=true, urlcolor=blue}

\begin{document}

\begin{center}
    {\LARGE \textbf{Yu Zhang}}\\[0.4em]
    GitHub: \href{https://github.com/CorradoZhang}{CorradoZhang} \,|\, E-mail: \href{mailto:Y.Zhang34@universityofgalway.ie}{Y.Zhang34@universityofgalway.ie} \,|\, LinkedIn: \href{https://www.linkedin.com/in/yu-zhang-ml-rl-rbt/}{yu-zhang-ml-rl-rbt}
\end{center}

\section*{Educational Background}

\textbf{University of Galway} \hfill Galway, Ireland\\
\textit{Ph.D.\ in Computer Science} \hfill Sept 2025 -- Sept 2029 (expected)

\textbf{University of Galway} \hfill Galway, Ireland\\
\textit{M.Sc.\ in Computer Science -- Artificial Intelligence (First Class Honours)} \hfill Sept 2024 -- Sept 2025

\textbf{Heilongjiang University of Science and Technology} \hfill Harbin, China\\
\textit{B.Eng.\ in Computer Science and Technology (Major Degree)} \hfill Oct 2020 -- Jun 2024\\
GPA: 3.55 (85.5\%)

\textit{B.Sc.\ in Mathematics and Applied Mathematics (Minor Degree)}\\
GPA: 3.75 (87.5\%)

\section*{Research Interests}

Reinforcement Learning, Robotics

\section*{Research Paper}

\begin{itemize}
    \item \textbf{Yu Zhang}, Karl Mason. ``D-SPEAR: Dual-Stream Prioritized Experience Adaptive Replay for Stable Reinforcement Learning in Robotic Manipulation.'' \textit{arXiv:2603.27346}, 2026. (Accepted at ICCRE 2026), [\href{https://arxiv.org/abs/2603.27346}{Link}]
\end{itemize}

\section*{Research \& Project Experience}

\textbf{D-SPEAR: Dual-Stream Prioritized Experience Adaptive Replay for Stable Reinforcement Learning in Robotic Manipulation} \hfill Oct 2025 -- Mar 2026\\
GitHub: \href{https://github.com/CorradoZhang/D-SPEAR-Code-and-Result}{https://github.com/CorradoZhang/D-SPEAR-Code-and-Result}
\begin{itemize}
    \item Designed and implemented D-SPEAR, a Soft Actor-Critic based reinforcement learning algorithm with a dual-stream prioritized replay mechanism, emphasizing high-TD samples for the critic, low-TD samples for the actor, and a shared anchor subset to stabilize updates across both branches.
    \item Implemented an adaptive replay coefficient $\lambda$ to adjust the balance between uniform and prioritized sampling based on the coefficient of variation of TD errors; additionally integrated candidate-ratio sampling and Huber critic loss to improve robustness and reward-scale adaptability.
    \item Built a robotic manipulation benchmark pipeline with robosuite and conducted 5-seed experiments on Lift and Door; measured by the mean return over the last 10 episodes, D-SPEAR achieved 305.4 on Lift, outperforming the best baseline TD3 (164.9) by 85.3\% and SAC (156.9) by 94.7\%.
    \item Completed the full research workflow including ablation studies, parallel multi-seed training, result aggregation, and publication-ready visualization; on Lift, removing the dual-stream design, low-TD actor sampling, or high-TD critic sampling reduced the mean return to 199.9, 179.0, and 184.2, corresponding to drops of about 34.6\%, 41.4\%, and 39.7\% from the full model (305.4), validating the effectiveness of the core design.
\end{itemize}

\textbf{Design of Crack Detection System for USTH Road Based on Image Processing} \hfill Harbin, China\\
\textit{Graduation Capstone Project for B.Eng.\ in Computer Science and Technology, Score: 93/100 (Rank: 1/158)}
\begin{itemize}
    \item Developed a Python-based road crack image collection system, enabling automated image capture both inside and outside the Heilongjiang University of Science and Technology campus.
    \item Designed an image processing system in MATLAB to process and display collected images, incorporating techniques such as histogram equalization and binary filtering.
    \item Implemented automatic crack severity classification and task assignment to relevant personnel.
    \item Built a web-based platform using Vue.js and Spring Boot, supporting CRUD operations for road crack data and improving usability for maintenance teams.
\end{itemize}

\section*{Extracurricular Activities \& Leadership}

\textbf{Heilongjiang University of Science and Technology AI Association} \hfill Sep 2021 -- Sep 2022\\
\textit{Technical Consultant, Training Group Leader}
\begin{itemize}
    \item Guided and mentored a team of 60+ junior members, focusing on practical skills in machine learning, data analysis, and model optimization.
    \item Provided technical oversight for team-led AI projects through code reviews, algorithm optimization, and project presentations.
    \item Fostered a collaborative learning environment by creating a knowledge-sharing platform, encouraging members to discuss challenges and share innovative solutions.
\end{itemize}

\section*{Awards \& Distinction}

\begin{itemize}
    \item Excellent Capstone Project (Rank: 1/158), Heilongjiang University of Science and Technology \hfill Jun 2024
    \item Second Prize, The Chinese Mathematics Competition (CMC), China Mathematical Society \hfill Mar 2023
    \item Second Prize, National English Competition for College Students (NECCS), International English Language Teachers Association \hfill Oct 2022
    \item Second Prize, MathorCup National College Mathematical Modeling Challenge Competition, Chinese Society of Optimization, Overall Planning and Economic Mathematics \hfill Oct 2022
    \item Third Prize, Lanqiao Cup National Software and Information Technology Professionals Competition, Talent Exchange Center Of Ministry Of Industry And Information Technology \hfill May 2022
    \item Third Prize, MathorCup National College Mathematical Modeling Challenge Competition -- Big Data Competition, Chinese Society of Optimization, Overall Planning and Economic Mathematics \hfill Mar 2022
    \item Third Prize, Northeast Three Provinces Mathematical Modeling League, Northeast Three Provinces Mathematical Modeling League Organizing Committee \hfill Aug 2021
\end{itemize}

\section*{Skills \& Language}

\textbf{Programming Skills:} Python (Advanced), C (Advanced), Java (Intermediate)\\
\textbf{Python Frameworks:} PyTorch, TensorFlow, Keras, Scikit-Learn, NumPy, Pandas, Matplotlib, Seaborn, SciPy\\
\textbf{Professional Applications:} MATLAB, LINGO, SPSS, CoppeliaSim\\
\textbf{Language:} English (Proficient, IELTS 6.5); Mandarin (Native)

\end{document}
